\begin{tikzpicture}

	%%% correlation of enhancer strength and gene expression data (without light) %%%
	\coordinate (gene expression) at (0, 0);
	
	\tikzset{
		leaves/.code = {
			\pgfmathsetlength{\templength}{-0.5\templength + 0.2 * rand * \templength}
			\pgfkeysalso{
				OkabeIto4,
				xshift = \templength
			}
		},
		other/.code = {
			\pgfmathsetlength{\templength}{0.5\templength + 0.2 * rand * \templength}
			\pgfkeysalso{
				black,
				xshift = \templength
			}
		},
		set box width/.code = {
			\pgfmathsetlength{\templength}{(\pgfkeysvalueof{/pgfplots/width} / 4) * 0.4}
		}
	}
	
	\def\tissue{\textbf{tissue:}\hspace*{-\groupplotsep}}
	\def\leaves{leaves}
	\def\other{other}

	\begin{pggfplot}[%
		at = {(gene expression)},
		total width = \threequartercolumnwidth,
		xlabel = species,
		ylabel = Pearson correlation coefficient,
		scatter styles = {tissue[empty legend],leaves,empty legend,other},
		enlarge x limits = 0,
		enlarge y limits = {lower = 0.1, upper = 0.1},
		xticklabels are csnames,
		title is csname,
		title color from name,
		xticklabel style = {font = \specieslongfalse},
		legend columns = 2,
		legend image post style = {xshift = -\templength},
		legend style = {name = legend, at = {(1, 1)}},
		skip facet = 1
	]{tamsACR_expression}
	
		\pggf_annotate(zeroline)
	
		\pggf_boxplot(
			color from column name,
			draw = black,
			forget plot,
			hide outliers
		)
	
		\pggf_scatter(
			style from column = tissue,
			mark size = 1.5,
			opacity = 0.75,
			set box width
		)
	
		\pggf_annotate(
			scatter legend*,
			only facet = 1
		)
	
	\end{pggfplot}
	
	
	%%% differentially expressed genes in light and dark %%%
	\coordinate (light specificity) at (gene expression -| \textwidth - \fourcolumnwidth, 0);
	
	\planticon[height = .5cm]{arabidopsis} at ($(light specificity) + (.6, -.25)$);
	
	\expandafter\def\csname dark_down\endcsname{downregulated}
	\expandafter\def\csname dark_up\endcsname{upregulated}
	
	\begin{pggfplot}[%
		at = {(light specificity)},
		total width = \fourcolumnwidth,
		xlabel = gene expression in dark,
		ylabel = linked ACR sequences (\%),
		ymin = 0,
		enlarge x limits = 0,
		enlarge y limits = {lower = 0, upper = .2},
		xticklabels are csnames,
		col title = light specificity,
		col title style = {left color = light!20, right color = dark!20, middle color = white},
		legend columns = 3,
		legend style = {anchor = north, at = {(.5, 1)}}
	]{tamsACR_dark_light_specificity}
	
		\pggf_bar(
			ybar,
			draw = black,
			style from column = {fill = group}
		)
	
	\end{pggfplot}
	
	
	%%% condition-/tissue-specificity of genes and corresponding ACRs %%%
	\coordinate[yshift = - \columnsep] (specificity) at (gene expression |- xlabel.south);
	
	\begin{pggfplot}[%
		at = {(specificity)},
		total width = \fourcolumnwidth,
		xlabel = tissue specificity,
		ylabel = condition specificity,
		enlarge y limits = {lower = 0.1, upper = 0.1},
		col title = specificity,
	]{tamsACR_specificity}
	
		\pggf_annotate(zeroline)
	
		\pggf_violin(
			fill = OkabeIto6,
			shade = left,
			shade color = white,
			cld = {anchor = south, text depth = 0pt}{max}
		)
	
	\end{pggfplot}
	
	
	%%% enriched GO terms %%%
	\coordinate (GO terms) at (specificity -| \textwidth - \threequartercolumnwidth, 0);
	
	\distance{col title.south}{col title.north}
	
	\begin{pggfplot}[%
		at = {(GO terms)},
		total width = \threequartercolumnwidth,
		xlabel = species,
		ylabel = GO term,
		ylabel add to width = 70pt,
		xticklabels are csnames,
		title is csname,
		title color from name,
		xticklabel style = {font = \specieslongfalse},
		y dir = reverse,
		enlargelimits = 0,
		colormap = {TF}{color(0) = (gray!20) color(1) = (OkabeIto4)},
		yticklabel style = {name = yticklabel},
		legend style = {anchor = east, at = {(0, 1)}, yshift = .5\ydistance - .5\pgflinewidth},
		legend columns = 3
	]{tamsACR_GO-terms}
	
		\pggf_heatmap(
			forget plot
		)
		
		\pggf_annotate(
			manual legend = {
				{\textbf{significant:}\hspace*{-\groupplotsep}} = {empty legend},
				yes = {area legend, fill = OkabeIto4},
				no = {area legend, fill = gray!20}
			},
			only facet = 1
		)
	
	\end{pggfplot}
	
	
	%%% subfigure labels %%%
	\subfiglabel{gene expression}
	\subfiglabel{light specificity}
	\subfiglabel{specificity}
	\subfiglabel{GO terms}
	
\end{tikzpicture}